\documentclass[main.tex]{subfiles}

% ==========================================
% Individual Lecture Title
% ==========================================

\ifSubfilesClassLoaded{%
    \title{Math 113: Field Extensions}
    \author{\defaultauthor}
    \date{November 17, 2026}
}{}

\begin{document}

\ifSubfilesClassLoaded{\maketitle}{}

% ==========================================
\section{November 17: Field Extensions}
% ==========================================

\subsection{Extension Fields}
Goal: Let $F$ be a field and $f(x) \in F[x]$ a nonconstant polynomial.

Show that there exists an \textbf{extension field} $E \supseteq F$ containing a zero $\alpha$ of $f(x)$
\subsection{Kronecker's Theorem}

\begin{theorem}[Kronecker's Theorem]
Let $F$ be a field and $f(x)\in F[x]$ a nonconstant polynomial.
Then there exists an extension field $E\supseteq F$ and an element
$\alpha\in E$ such that $f(\alpha)=0$.
\end{theorem}

\begin{proof}
Since $F$ is a field, we can factor $f(x)$ into irreducible
polynomials over $F$. Let $p(x)$ be one of these irreducible factors.
It suffices to find a root $\alpha$ of $p(x)$, since then
$f(\alpha)=0$ as well.

Because $p(x)$ is irreducible in $F[x]$, the ideal $\langle p(x)\rangle$
is maximal. Therefore
\[
E\defeq F[x]/\langle p(x)\rangle
\]
is a field.

We want an injective ring homomorphism $\psi:F\to E$. Define
\[
\psi(a)=a+\langle p(x)\rangle.
\]

This is a ring homomorphism because it is the composition
\[
F\longrightarrow F[x]\longrightarrow F[x]/\langle p(x)\rangle
\]
of the inclusion of constants and the quotient map.

To prove injectivity, suppose $\psi(a)=\psi(b)$. Then
\[
a+\langle p(x)\rangle=b+\langle p(x)\rangle,
\]
so $a-b\in\langle p(x)\rangle$.

Since $a-b$ is constant while $p(x)$ has positive degree, this
is possible only if $a-b=0$.

Thus $a=b$ and $\psi$ is injective. We may therefore identify
$F$ with $\im(\psi)\subseteq E$.

Now let
\[
\alpha=x+\langle p(x)\rangle\in E.
\]

Write $p(x)=a_0+a_1x+\cdots+a_nx^n$, with $a_i\in F\subseteq E$.

Then
\begin{align*}
p(\alpha)
 &=a_0+a_1\alpha+\cdots+a_n\alpha^n\\
 &=\bigl(a_0+a_1x+\cdots+a_nx^n\bigr)+\langle p(x)\rangle\\
 &=p(x)+\langle p(x)\rangle\\
 &=0_E.
\end{align*}

Therefore $\alpha$ is a root of $p$, and hence of $f$.
\end{proof}


% ==========================================
% Examples
% ==========================================

\begin{example}
Let $F=\R$ and $f(x)=x^2+1$.

This polynomial is irreducible over $\R$, so
\[
E=\R[x]/\langle x^2+1\rangle
\]
is a field.

The element $\alpha=x+\langle x^2+1\rangle$ is a root of $f$
in $E$, because
\[
\alpha^2+1=(x^2+1)+\langle x^2+1\rangle=0_E.
\]

In fact, $E\cong\C$ via the map sending
$x+\langle x^2+1\rangle$ to $i$.
\end{example}


\begin{example}
Let $F=\Q$ and
\[
f(x)=x^4-5x^2+6=(x^2-2)(x^2-3).
\]

We may choose either irreducible factor
$p(x)=x^2-2$ or $p(x)=x^2-3$ to construct a field
extension containing a root of $f$.
\end{example}


% ==========================================
\subsection{Algebraic and Transcendental Elements}
% ==========================================

We can classify elements of a field extension $E/F$ into two kinds.

\begin{definition}
Let $E/F$ be a field extension. An element $\alpha\in E$ is
\emph{algebraic over $F$} if $f(\alpha)=0$ for some nonzero
$f(x)\in F[x]$.

Otherwise, $\alpha$ is \emph{transcendental over $F$}.
\end{definition}


\begin{example}
The element $i\in\C$ is algebraic over $\R$, since
$i^2+1=0$.

It is also algebraic over $\Q$, for the same reason.
\end{example}


\begin{example}
The element $\sqrt{2}\in\R$ is algebraic over $\Q$, since
it is a root of $x^2-2\in\Q[x]$.
\end{example}


\begin{example}
The number $\pi\in\R$ is transcendental over $\Q$.
\end{example}


\begin{example}
The number $\pi$ is algebraic over $\Q(\pi)$, the smallest
subfield of $\R$ containing $\Q$ and $\pi$, since it is a root
of $x-\pi\in\Q(\pi)[x]$.
\end{example}


\begin{example}
Let $\alpha=\sqrt{1+\sqrt{3}}\in\R$.

Then $\alpha$ is algebraic over $\Q$, since
\[
\begin{aligned}
\alpha^2&=1+\sqrt{3}\\
(\alpha^2-1)^2&=3\\
(\alpha^2-1)^2-3&=0.
\end{aligned}
\]

Thus $\alpha$ is a root of the nonzero polynomial
$(x^2-1)^2-3\in\Q[x]$.
\end{example}


An element of $\C$ that is algebraic over $\Q$ is called an
\emph{algebraic number}.


% ==========================================
% Transcendental Elements and Homomorphisms
% ==========================================

\begin{theorem}
Let $E/F$ be a field extension and $\alpha\in E$.

Then $\alpha$ is transcendental over $F$ if and only if the
evaluation homomorphism
\[
\phi_\alpha:F[x]\longrightarrow E,
\qquad f(x)\longmapsto f(\alpha)
\]
is injective.
\end{theorem}


\begin{proof}
\[
\begin{aligned}
\alpha\text{ is transcendental over }F
&\iff f(\alpha)\ne0
\quad\text{for every nonzero }f(x)\in F[x]\\
&\iff \phi_\alpha(f(x))\ne0
\quad\text{for every nonzero }f(x)\in F[x]\\
&\iff \ker(\phi_\alpha)=\{0\}\\
&\iff \phi_\alpha\text{ is injective}.
\end{aligned}
\]
\end{proof}


% ==========================================
% Minimal Polynomials
% ==========================================

\begin{theorem}
Let $E/F$ be a field extension and let $\alpha\in E$ be
algebraic over $F$.

There is an irreducible polynomial $p(x)\in F[x]$ such that
$p(\alpha)=0$.

This polynomial is unique up to multiplication by a nonzero
constant in $F$, and has the smallest degree among all nonzero
polynomials in $F[x]$ having $\alpha$ as a root.

Moreover, for every $f(x)\in F[x]$,
\[
f(\alpha)=0\quad\Longrightarrow\quad p(x)\mid f(x).
\]
\end{theorem}


\end{document}